% User-provided Irodov PDF pp.70, printed p.61;
% Section 1.8 final problems 381--388.
% Source answer p.295 has a unit error in problem 386: 1.43e3 GeV,
% not 1.43e3 MeV, for T=25 GeV per proton.
\exproblem{381}{61}
一个相对论性粒子在惯性系 $K$ 中沿
$x$ 轴运动，动量为 $p_x$、总能量为 $E$。
证明在相对 $K$ 沿 $+x$ 方向以恒速 $V$
运动的另一惯性系 $K'$ 中，
粒子动量和总能量变换为
\[
p'_x=\frac{p_x-EV/c^2}{\sqrt{1-V^2/c^2}},
\qquad
E'=\frac{E-Vp_x}{\sqrt{1-V^2/c^2}}.
\]
\exsolution
粒子在自身世界线上的固有时间
$\mathrm d\tau$ 为洛伦兹标量。
四动量定义为
\[
P^\mu=m_0\frac{\mathrm dx^\mu}{\mathrm d\tau},
\quad
x^\mu=(ct,x,y,z),
\quad
P^\mu=(E/c,p_x,p_y,p_z).
\]
坐标的洛伦兹变换为
\[
ct'=\gamma_V(ct-\beta x),\qquad
x'=\gamma_V(x-Vt),
\quad
\gamma_V=(1-\beta^2)^{-1/2},\quad\beta=V/c.
\]
由于分母 $\mathrm d\tau$ 不变，
对坐标变换沿世界线求导可知四动量
具有完全相同的线性变换律：
\[
P'^0=\gamma_V(P^0-\beta P^1),\qquad
P'^1=\gamma_V(P^1-\beta P^0).
\]
代回物理分量：
\[
\boxed{E'=\gamma_V(E-Vp_x)},\qquad
\boxed{p'_x=\gamma_V(p_x-VE/c^2)}.
\]
\excheck 四动量变换同时适用于有质量粒子和光子；
不能将能量视为标量而使其在各惯性系不变。

\exproblem{382}{61}
一光子在惯性系 $K$ 中的能量为
$\varepsilon$。
惯性系 $K'$ 沿光子传播方向以速度 $V$
相对 $K$ 运动。
用上一题四动量变换求光子在 $K'$ 中的能量
$\varepsilon'$；并求使
$\varepsilon'=\varepsilon/2$ 的 $V$。
\exsolution
光子沿正向传播，动量
$p_x=\varepsilon/c$。
由四动量能量变换，
\[
\varepsilon'=\gamma_V(\varepsilon-Vp_x)
=\gamma_V\varepsilon(1-\beta)
=\boxed{\varepsilon\sqrt{\frac{1-\beta}{1+\beta}}},
\qquad \beta=V/c.
\]
令 $\varepsilon'/\varepsilon=1/2$，
平方得
$(1-\beta)/(1+\beta)=1/4$，
故 $4-4\beta=1+\beta$，
\[
\boxed{V=\frac35c}.
\]
\excheck 与光子同向运动的观察者测得红移；
若 $V$ 逆光子方向，变换应得到蓝移而不是这一公式。

\exproblem{383}{61}
证明对于有静质量的粒子，组合
$E^2-p^2c^2$ 是所有惯性系中的不变量，
并给出这一不变量的数值。
\exsolution
在粒子静止系有
$E_0=m_0c^2$、$\bm p_0=\bm0$，
故
$E_0^2-p_0^2c^2=m_0^2c^4$。
由上一题的四动量洛伦兹变换，
\[
\begin{aligned}
E'^2-c^2p_x'^2
&=\gamma_V^2[(E-Vp_x)^2-
c^2(p_x-VE/c^2)^2]\\
&=\gamma_V^2(1-\beta^2)(E^2-c^2p_x^2)
=E^2-c^2p_x^2.
\end{aligned}
\]
垂直方向动量不变，因此一般三维情况下
$E'^2-c^2|\bm p'|^2=E^2-c^2|\bm p|^2$。
结论为
\[
\boxed{E^2-p^2c^2=m_0^2c^4}.
\]
\excheck 静止质量 $m_0$ 与惯性系无关，
而粒子的总能量、动量分别随参考系改变。

\exproblem{384}{61}
一颗静质量为 $m_0$ 的中子动能为
$T=2m_0c^2$，撞向另一颗实验室静止的中子。
求：(a) 两中子在质心系的总动能
$\widetilde T$ 与每颗中子的动量模
$\widetilde p$；
(b) 这个质心系相对实验室的速度。
\exsolution
实验室中，入射中子总能量
$E_1=m_0c^2+T=3m_0c^2$，
靶中子总能量 $E_2=m_0c^2$。
故系统总能量
$E_{\rm tot}=4m_0c^2$。
入射动量模由
$p_1c=\sqrt{E_1^2-m_0^2c^4}
=2\sqrt2\,m_0c^2$。
系统四动量不变量为
\[
E_{\rm cm}^2
=E_{\rm tot}^2-p_{\rm tot}^2c^2
=(16-8)m_0^2c^4=8m_0^2c^4.
\]
因此
$E_{\rm cm}=2\sqrt2\,m_0c^2$。
去掉两颗静止能量得质心系总动能：
\[
\boxed{\widetilde T
=2(\sqrt2-1)m_0c^2
\simeq777\,\mathrm{MeV}}.
\]
两中子在质心系总动量为零且质量相等，
每颗总能量为
$E_{\rm cm}/2=\sqrt2 m_0c^2$，
因此
\[
\boxed{\widetilde p=
\frac1c\sqrt{2m_0^2c^4-m_0^2c^4}
=m_0c\simeq940\,\mathrm{MeV}/c}.
\]
质心速度可由实验室总能量与总动量确定：
\[
\boxed{V_{\rm cm}
=\frac{p_{\rm tot}c^2}{E_{\rm tot}}
=\frac{c}{\sqrt2}
\simeq2.12\times10^8\,\mathrm{m/s}}.
\]
\excheck 必须使用系统的\emph{四动量不变量}
而非只在实验室中对两颗中子动能做算术平均。

\exproblem{385}{61}
静质量为 $m_0$、动能为 $T$ 的粒子
撞上一颗实验室静止、具有相同静质量的粒子，
两者完全结合为一个复合粒子。
求复合粒子的静质量 $M_0$
及它相对于实验室的速度 $V$。
\exsolution
碰前系统总能量与动量为
\[
E_{\rm tot}=(2m_0c^2+T),\qquad
p_{\rm tot}c=\sqrt{T(T+2m_0c^2)}.
\]
碰后复合粒子的静质量由四动量不变量定义：
\[
M_0^2c^4=E_{\rm tot}^2-p_{\rm tot}^2c^2
=(2m_0c^2+T)^2-T(T+2m_0c^2).
\]
展开后
$M_0^2c^4=4m_0^2c^4+2m_0c^2T$，
\[
\boxed{M_0=\sqrt{2m_0(2m_0+T/c^2)}}.
\]
复合粒子速度由
$\bm V=\bm p_{\rm tot}c^2/E_{\rm tot}$，
故
\[
\boxed{V=c\sqrt{\frac{T}{T+2m_0c^2}}}.
\]
\excheck $T>0$ 时 $M_0>2m_0$，
多出的静质量体现碰撞耗散的内部能量；
若把它直接写成 $2m_0$，违反总能量守恒。

\exproblem{386}{61}
问一个高能质子撞向实验室中\emph{静止}的另一质子时，
入射质子动能 $T_{\rm beam}$ 至少需要多高，
才能使该过程的\emph{质心系总动能}
等于两束相向质子、每束动能
$T=25.0\,\mathrm{GeV}$ 时的质心系总动能？
\exsolution
两质子相向且各自动能为 $T$ 时，
质心系就是两束的对称实验室系。
质心系总能量为
$E_{\rm cm}=2(m_pc^2+T)$。
对应的固定靶碰撞中，
实验室总能量
$E_{\rm lab}=2m_pc^2+T_{\rm beam}$，
入射动量满足
$p_{\rm lab}^2c^2
=T_{\rm beam}(T_{\rm beam}+2m_pc^2)$。
四动量不变量为
\[
s=E_{\rm lab}^2-p_{\rm lab}^2c^2
=4m_p^2c^4+2m_pc^2T_{\rm beam}.
\]
令它等于对撞束的
$s=4(m_pc^2+T)^2$，
得到
\[
2m_pc^2T_{\rm beam}
=4[(m_pc^2+T)^2-m_p^2c^4]
=4T(T+2m_pc^2).
\]
所以
\[
\boxed{T_{\rm beam}=
\frac{2T(T+2m_pc^2)}{m_pc^2}}.
\]
用质子静能 $m_pc^2\simeq0.938\,\mathrm{GeV}$：
\[
\boxed{T_{\rm beam}
\simeq\frac{2(25.0)(25.0+1.876)}{0.938}
\,\mathrm{GeV}\simeq1.43\times10^3\,\mathrm{GeV}
=1.43\,\mathrm{TeV}}.
\]
\excheck \textbf{本版答案单位冲突：}
书末公式与原题 $T=25\,\mathrm{GeV}$ 相容，
但数值旁印 $1.43\times10^3\,\mathrm{MeV}$；
按单位计算应为
$1.43\times10^3\,\mathrm{GeV}$，
相差1000倍。本稿保留物理量纲正确的值。

\exproblem{387}{61}
一颗静质量为 $m_0$ 的粒子
原来在实验室静止，
自发衰变为静质量分别为
$m_1,m_2,m_3$ 的三颗粒子。
求粒子1可能获得的最大总能量。
\exsolution
初态总四动量为
$(m_0c,\bm0)$。
将粒子2与3组成一个复合末态系统，
其不变质量记为 $M_{23}$。
总静质量下界
$M_{23}\ge m_2+m_3$；
仅当粒子2、3相对静止、
具有同一速度时取等。
把衰变视为
$m_0\to m_1+M_{23}$ 的二体末态。
在初始粒子静止系，由
$P_{23}=P_0-P_1$ 和
$P_{23}^2=M_{23}^2c^2$，
\[
M_{23}^2c^4
=(m_0c^2-E_1)^2-p_1^2c^2
=m_0^2c^4+m_1^2c^4-2m_0c^2E_1.
\]
于是
\[
E_1=
\frac{m_0^2+m_1^2-M_{23}^2}{2m_0}c^2.
\]
为使 $E_1$ 最大，令 $M_{23}=m_2+m_3$：
\[
\boxed{E_{1,\max}=
\frac{m_0^2+m_1^2-(m_2+m_3)^2}{2m_0}c^2}.
\]
\excheck 该值可达到的必要条件是
$m_0\ge m_1+m_2+m_3$；
要最大化单个粒子的能量，
\emph{其余两粒子须作为一个共动整体}
使不变质量最小。

\exproblem{388}{61}
一枚相对论火箭沿直线喷出气体，
气体相对火箭的速度大小固定为
$u$ 且满足 $u\ll c$。
火箭的静质量由初始 $m_0$
逐渐减到 $m$，无外力。
求火箭相对初始静止系的速度 $v(m)$。
\exsolution
每次考虑一个火箭的\emph{瞬时共动惯性系}。
以沿推进方向的火箭速度增量为正，
喷气速度为 $-u$。
在非相对论喷流近似下，
微元动量守恒得到
$m\,\mathrm dv_{\rm com}=-u\,\mathrm dm$，
其中 $\mathrm dm<0$。
为使不同瞬时共动系中的速度增量相加，
引入快度
$\chi=\operatorname{artanh}(v/c)$。
快度差与瞬时共动速度增量满足
$\mathrm d\chi=\mathrm dv_{\rm com}/c$，
因此
\[
\mathrm d\chi=-\frac{u}{c}\frac{\mathrm dm}{m}.
\]
从 $m_0$、初速零积分到 $m$：
\[
\chi=\frac uc\ln\frac{m_0}{m}.
\]
最后 $v/c=\tanh\chi$，
利用
$\tanh z=(1-e^{-2z})/(1+e^{-2z})$：
\[
\boxed{\frac vc=
\tanh\!\left(\frac uc\ln\frac{m_0}{m}\right)
=\frac{1-(m/m_0)^{2u/c}}
{1+(m/m_0)^{2u/c}}}.
\]
\excheck 当 $u/c$ 很小且累计速度也远小于 $c$，
$\tanh z\simeq z$，
退化为通常的火箭方程
$v\simeq u\ln(m_0/m)$；
在严格相对论模型中 $v<c$ 恒成立。
